This paper addresses a specific gap in physical climate risk assessment: the absence of duration dependence in standard depth-damage functions used by CLIMADA, JRC, and regulatory stress tests. Using 295{,}637 FEMA National Flood Insurance Program claims after conservative quality filters and depth capping, we calibrate a duration-dependent extension of the Richards damage function and demonstrate its application to corporate credit risk.

The main findings are:

\begin{enumerate}
    \item \textbf{Duration increases flood damage by up to 2.6 times.} At 1~m depth, the calibrated model predicts DR = 0.25 for instantaneous flooding, rising to 0.66 for 288-hour inundation. This effect operates primarily through inflection point shift ($\lambda_a = 0.49$); the duration-modifying coefficients $\lambda_Q \approx 0$ and $\delta_\nu \approx 0$ are not significantly non-zero in the three-event calibration, suggesting that shape modification may be secondary to inflection shift for these hurricane events, though this conclusion is limited by the small number of distinct $\tau$ values available. Standard JRC/CLIMADA functions, which lack duration, underpredict prolonged flood damage by about 15 percentage points (about 50\% MAE reduction with our model).

    \item \textbf{Sector adjustment $\kappa$ is calibrated from 82{,}650 NFIP non-residential claims.} Grouping by building stories and insurance coverage, we obtain an empirical lookup table for building-level vulnerability. A coverage-gradient extrapolation to corporate scale ($\rho_{\text{gradient}} = 0.25$) bridges NFIP-insured commercial buildings to large corporate facilities. Out-of-sample validation on 24 companies across six flood events yields median ratio $0.90\times$ (23/24 within 0.4--2.5$\times$), with no corporate data entering the calibration.

    \item \textbf{Indirect losses are material for event-scale screening.} CapEx-only frameworks underestimate total loss when operational disruption is significant. Applied to ADM's 12 flood-exposed facilities, the calibrated pipeline produces facility-level loss estimates with duration effects, CapEx/OpEx decomposition, and insurance offsets.
\end{enumerate}

The calibrated parameters $(B = 5.65, Q_0 = 3.45, a_0 = 0.54, \nu_0 = 2.61, \lambda_a = 0.49)$ are directly usable by practitioners and can serve as inputs to any climate risk platform that currently employs static depth-damage curves.

Several additional limitations warrant acknowledgement.

\paragraph{Duration assignment} The flood duration $\tau$ is assigned per event (Sandy $\sim$36~h, Harvey $\sim$288~h, Katrina $\sim$120~h), not per individual claim or property. Importantly, event-level $\tau$ is a proxy for local inundation duration: actual water-contact time varies substantially within a single flood event, driven by local topography, drainage infrastructure, and distance from the flood source. Our calibration captures \textit{inter-event} duration variation across events spanning an 8 times range in $\tau$, but cannot resolve \textit{intra-event} heterogeneity. Although the current OpenFEMA v2 schema includes a claim-level \texttt{floodWaterDuration} field, a targeted audit of 50{,}000 raw claims per event found effectively no positive duration values in the Sandy, Harvey, and Katrina extracts, so the field was not informative for direct calibration. We therefore rely on hydrological literature and FEMA flood extent reports for event-level estimates. Future work with georeferenced flood recession data could enable spatially varying duration assignment within single events.

\paragraph{Individual claim accuracy} The damage function predicts \textit{average} damage ratios at a given depth and duration. Individual claims exhibit substantial variance (MAE $\approx 0.25$, R$^2 \approx 0.05$ at claim level), reflecting factors not captured by depth alone. Feature enrichment reduces MAE by only 11.7\%, with an irreducible noise floor of $\sim$0.22. The model is best suited to event-scale and portfolio-level applications, where individual errors average out and the strong ZIP-level correlation ($\rho = 0.949$) supports geographic risk ranking.

\paragraph{Hazard characterisation} This paper uses inundation depth and duration as the primary intensity measures. Flow velocity, which can contribute substantially to structural damage in riverine and coastal flood events, is not captured. For events where velocity is a major damage driver, the depth-only Richards formulation will underestimate damage, and an extension incorporating velocity (as in~\cite{merz2010review}) would be needed.

\paragraph{Geographic scope} The damage function is calibrated exclusively on US flood insurance claims. Transferability to other countries -- with different construction standards (e.g., brick masonry in Europe vs.\ wood frame in the US), flood mechanisms (riverine vs.\ pluvial vs.\ coastal), and insurance practices -- requires separate validation.

Promising directions for future work include: (i)~moving from event-level to spatially resolved duration assignment within single flood events; (ii)~extending the bootstrap treatment to hierarchical and spatially explicit uncertainty models; (iii)~validating transferability to non-US construction standards; and (iv)~coupling the calibrated damage function to richer downstream disruption and recovery models.
